\begin{section}{Erweiterungen für MOROS}
\label{sec:erweiterungen}

Um ein interaktives, verteiltes Spiel auf MOROS zu realisieren, waren mehrere Systemerweiterungen nötig:
(1) eine PS/2-Maus\-unterstützung inklusive IRQ-Handler und Event-Queue,
(2) ein einfacher Framebuffer mit Zeichenprimitiven und Textausgabe,
(3) Ergänzungen und Bugfixes im UDP-Stack (\texttt{listen}-Methode und zuverlässigeres \texttt{write}),
sowie (4) Kernel-nahe Debug-Hilfen (\texttt{kprint!}/\texttt{kprintln!}).
Diese Bausteine bilden das Fundament für Eingabe, Ausgabe und Netzwerk des Spiels.
Eine ähnliche Übersicht wurde bereits in unserem Vortrag gegeben.

\begin{subsection}{Framebuffer: Grafikprimitive und Text}
\label{subsec:fb}
Der Framebuffer ist als rein userspace-seitig beschriebenes Bytefeld implementiert, das über \texttt{flush()} atomar in das vom Treiber beobachtete Gerätedatei-Interface geschrieben wird.
Die Struktur kapselt Auflösung, Farbtiefe (hier: 8 Bit/Pixel) und \emph{pitch} (Bytes pro Zeile):
\begin{verbatim}
pub struct Framebuffer {
    width, height, color_depth, pitch, file_path, internal_buffer, buffer_size
}
\end{verbatim}

\paragraph{Speicherlayout und Flush.}
\texttt{pitch = width * (color\_depth / 8)}; der interne Puffer hat exakt \texttt{height * pitch} Bytes.
\texttt{flush()} persistiert den gesamten Puffer:
\begin{verbatim}
pub fn flush(&self) {
    write(self.file_path, &self.internal_buffer).expect("...");
}
\end{verbatim}
Damit ist die Ausgabe wohldefiniert (kein „tearing“ durch Teilwrites).

\paragraph{Primitive: Pixel, Linie, Rechteck, Kreis.}
\texttt{draw\_pixel} prüft Grenzen und schreibt ein Byte (Farbindex) an \texttt{y * pitch + x}.
\texttt{draw\_line} nutzt den Bresenham-Algorithmus (integerbasiert, verzicht auf Floating-Point im Hotpath),
\texttt{draw\_rectangle} füllt O(\,Breite\,$\times$\,Höhe\,) Pixel,
\texttt{draw\_circle} implementiert die Midpoint-Variante; bei \texttt{filled=true} werden horizontale Spans gezeichnet:
\begin{verbatim}
self.draw_line(cx - x, cy + y, cx + x, cy + y, color);
... // vier Symmetrieachsen
\end{verbatim}

\paragraph{Textausgabe mit Bitmap-Font und Skalierung.}
Die Funktion \texttt{draw\_text} iteriert über \texttt{font.data} (8 Spalten pro Zeichenzeile) und setzt Pixel dort,
wo das Bitmuster 1 ist. Das optionale \texttt{scale} vergrößert Abstände und Pixelpositionen:
\begin{verbatim}
if (row >> (7 - dx)) & 1 == 1 {
    let px = (cursor_x + dx as f32 * scale) as isize;
    let py = (y as f32 + dy as f32 * scale) as isize;
    self.draw_pixel(px, py, color);
}
\end{verbatim}
\textbf{Hinweis (aktueller Stand):} Skalierung $>1$ repliziert Pixel nicht flächig (kein „Block-Fill“),
sondern setzt versetzte Einzelpixel; sehr kleine Skalen (<1) können Lücken erzeugen.
Für UI-Text im Spiel (Stats links, Koordinatenraster) genügt das, eine echte „Nearest-Neighbor“-Skalierung wäre jedoch eine sinnvolle Erweiterung.

\paragraph{Komposition: \texttt{copy\_from} und \texttt{copy\_from\_at}.}
Für HUD/Offscreen-Buffering können Inhalte kopiert werden.
\texttt{copy\_from\_at} iteriert zeilenweise und kopiert Byte-für-Byte;
bei 8-Bit-Farbtiefe ist das korrekt, für höhere Farbtiefen müsste die Schleife auf \texttt{bytes\_per\_pixel} erweitert werden (derzeit nicht nötig).
\end{subsection}

\begin{subsection}{PS/2-Maus: IRQ-Handler und Event-Queue}
\label{subsec:mouse}
Wir initialisieren den PS/2-Controller, aktivieren den zweiten Port und konfigurieren IRQ12.
Die Initialisierung sendet klassische PS/2-Kommandos: \texttt{0xA8} (Enable second PS/2), \texttt{0x20}/\texttt{0x60} (Config read/write),
\texttt{0xD4} (Write to mouse) und \texttt{0xF4} (Enable data reporting).
Eben diese Schritte sowie die Nutzung einer Queue wurden im Vortrag als Voraussetzung genannt.

\paragraph{Interrupt-Pfad.}
Der IRQ-Handler ruft \texttt{internal\_irq\_handler()}, das genau ein Dreibyte-Paket (\texttt{[status, dx, dy]})
aus den Ports \texttt{0x64}/\texttt{0x60} liest, prüft \emph{Output Buffer Full} und \emph{Always-1}-Bit,
und liefert ein \texttt{MouseEvent}:
\begin{verbatim}
Some(MouseEvent {
    x_movement: mouse_data[1] as i8,
    y_movement: -(mouse_data[2] as i8), // y-Achse invertiert
    buttons: mouse_data[0] & 0x07,
})
\end{verbatim}

\paragraph{Lockfreie Event-Queue.}
Zur Übergabe in den Spiel-Loop wird ein MPMC-Queue (bounded SCQ, Kapazität 128) verwendet:
\begin{verbatim}
let (receiver, sender) = mpmc::bounded::scq::queue(128);
sender.try_enqueue(event).ok();
\end{verbatim}
Der Client ruft \texttt{get\_last\_event()} (nicht-blockierend) oder \texttt{wait\_for\_event()} (blockierend) auf.
Die bounded Queue bietet deterministisches Verhalten ohne Heap-Wachstum und vermeidet Spinlocks im Hotpath.

\paragraph{Fehlertoleranz.}
Ist der Receiver geschlossen, wird defensiv abgebrochen (\texttt{panic!});
im produktiven Kernel wäre hier ein definierter Fallback wünschenswert.
Für das Seminar-Setup genügt das strikte Verhalten.
\end{subsection}

\begin{subsection}{UDP-Erweiterungen und Bugfixes}
\label{subsec:udp}
Für das verteilte Spiel benötigen wir zuverlässige, eng getaktete UDP-Send/Receive-Zyklen.
In MOROS ergänzen wir dazu \texttt{listen(port)} und korrigieren die Semantik von \texttt{write()},
so dass nach einem Send nicht „zu früh“ aus der Poll-Schleife ausgestiegen wird—ein Problem, das im Vortrag als Fehlerbild (Sekunden-lange Sendepausen) beschrieben wurde.

\paragraph{\texttt{listen(port)}.}
Die Funktion bindet den Socket, sobald dieser \emph{nicht} offen ist, und \emph{pollt} die Netzwerkschicht deterministisch:
\begin{verbatim}
let local = IpListenEndpoint::from(port);
socket.bind(local).map_err(|_| ())?;
\end{verbatim}
Zwischenpolling über \texttt{iface.poll\_delay(...)} und \texttt{halt()} verhindert Busy-Waiting,
ein Timeout schützt vor Hängern (hier 5\,s).

\paragraph{\texttt{write()}: zuverlässiger Sendeabschluss.}
\emph{Vorher} wurde nach dem ersten \texttt{send\_slice} ein Flag \texttt{sent} gesetzt und erst \emph{im nächsten} Poll-Durchlauf verlassen,
wodurch unter Last Sends „liegenbleiben“ konnten.
\emph{Nachher} pollt die Funktion \emph{sofort} erneut und bricht dann sauber ab:
\begin{verbatim}
if socket.can_send() {
    socket.send_slice(buf, endpoint)?;
    iface.poll(...); // sofortiger Fortschritt
    break;
}
\end{verbatim}
Damit werden Sendefenster zügig geleert; Timeouts und \texttt{poll\_delay} bleiben als Backpressure-Mechanismen erhalten.
\end{subsection}

\begin{subsection}{Debugging-Hilfen}
\label{subsec:debug}
Zur Diagnose im Kernelmodus verwenden wir \texttt{kprint!} und \texttt{kprintln!} (aus TOSR übernommen),
um frühe Boot- und Interruptpfade zu beobachten.
Im Vortrag sind diese Hilfen als separate Erweiterung aufgeführt; sie waren während der Netzwerk- und PS/2-Entwicklung essenziell.
\end{subsection}

\begin{subsection}{Zusammenfassung und Grenzen}
Die beschriebenen Bausteine liefern genau das, was das Spiel benötigt:
\emph{Frame\-buffer}-Zeichenprimitiven für Map/Spieler/Bullets,
\emph{Maus}-Events für Zielen/Schießen,
sowie \emph{UDP} für die State-Synchronisation.
Aktuelle Grenzen sind:
(1) 8-Bit-Farbpalette (bewusst, aber limitiert),
(2) einfache Textskalierung ohne Block-Fill,
(3) defensive \texttt{panic!}-Strategie bei Queue-Fehlern,
(4) UDP ohne Fragmentierungsstrategien (vgl. Analysekapitel).
Diese Punkte sind bewusst gewählte Vereinfachungen zugunsten eines stabilen, verständlichen Systems und decken sich mit der im Vortrag skizzierten Pragmatik.

\end{subsection}
\end{section}